% from ega4-17-fr.tex
\oldpage[IV-4]{58}
$Z$-morphisme $u_0:Z'_0\to X$; l'hypothèse sur $g$ entraîne qu'il existe un $Z$-morphisme $v:Z'\to Y$ tel que $f\circ u_0=v\circ j$ (où $j:Z'_0\to Z'$ est l'injection canonique); l'hypothèse sur $f$ entraîne alors qu'il existe un morphisme $u:Z'\to X$ tel que $f\circ u=v$ et $u\circ j=u_0$, donc $(g\circ f)\circ u$ est égal au morphisme donné $Z'\to Z$ et $u\circ j=u_0$, ce qui prouve que $g\circ f$ est formellement lisse; on raisonne de même lorsque l'on suppose $f$ et $g$ formellement non ramifiés.

Enfin, pour démontrer (iii), posons $X'=X_{(S')}$, $Y'=Y_{(S')}$, $f'=f_{(S')}$; considérons un schéma affine $Y''$, un sous-schéma fermé $Y''_0$ de $Y''$ défini par un Idéal nilpotent et un morphisme $g:Y''\to Y'$ faisant de $Y''$ un $Y'$-préschéma; on sait alors (I, 3.3.8) que $\Hom_{Y'}(Y'',X')$ s'identifie canoniquement à $\Hom_Y(Y'',X)$ et $\Hom_{Y'}(Y''_0,X')$ à $\Hom_Y(Y''_0,X)$, et la conclusion résulte alors immédiatement de la définition (17.1.1).

On notera qu'une \emph{immersion fermée} n'est pas nécessairement un morphisme formellement lisse.

\begin{proposition}[17.1.4]
\label{IV.17.1.4}
Soient $f:X\to Y$, $g:Y\to Z$ deux morphismes, et supposons $g$ formellement non ramifié. Alors, si $g\circ f$ est formellement lisse (resp. formellement étale), il en est de même de $f$.
\end{proposition}

En effet, soient $Y'$ un schéma affine, $Y'_0$ un sous-schéma fermé de $Y'$ défini par un Idéal nilpotent, $h:Y'\to Y$ un morphisme, $j:Y'_0\to Y'$ l'injection canonique, $u_0:Y'_0\to X$ un $Y$-morphisme, donc tel que $f\circ u_0=h\circ j$. Supposons $g\circ f$ formellement lisse; alors il existe un morphisme $u:Y'\to X$ tel que $u\circ j=u_0$ et $(g\circ f)\circ u=g\circ h$. Mais ces relations entraînent que $f\circ u$ et $h$ sont deux $Z$-morphismes de $Y'$ dans $Y$ tels que $(f\circ u)\circ j=h\circ j$; en vertu de l'hypothèse que $g$ est formellement non ramifié, on en tire que $f\circ u=h$, autrement dit $u$ est un $Y$-morphisme; donc $f$ est formellement lisse. Compte tenu de (17.1.3, (v)), cela démontre la proposition.

\begin{corollary}[17.1.5]
\label{IV.17.1.5}
Supposons $g$ formellement étale; alors, pour que $g\circ f$ soit formellement lisse (resp. formellement non ramifié, resp. formellement étale), il faut et il suffit que $f$ le soit.
\end{corollary}

Cela résulte de (17.1.4) et de (17.1.3, (ii) et (v)).

\begin{proposition}[17.1.6]
\label{IV.17.1.6}
Soit $f:X\to Y$ un morphisme de préschémas.
\begin{enumerate}
  \item[(i)] Soit $(U_\alpha)$ un recouvrement ouvert de $X$ et, pour tout $\alpha$, soit $i_\alpha:U_\alpha\to X$ l'injection canonique. Pour que $f$ soit formellement lisse (resp. formellement non ramifié, resp. formellement étale), il faut et il suffit que chacun des morphismes $f\circ i_\alpha$ le soit.
  \item[(ii)] Soit $(V_\lambda)$ un recouvrement ouvert de $Y$. Pour que $f$ soit formellement lisse (resp. formellement non ramifié, resp. formellement étale), il faut et il suffit que chacune des restrictions $f^{-1}(V_\lambda)\to V_\lambda$ de $f$ le soit.
\end{enumerate}
\end{proposition}

Notons d'abord que (ii) est conséquence de (i): en effet, si $j_\lambda:V_\lambda\to Y$ et $i_\lambda:f^{-1}(V_\lambda)\to X$ sont les injections canoniques, la restriction $f_\lambda:f^{-1}(V_\lambda)\to V_\lambda$ de $f$ est telle que $j_\lambda\circ f_\lambda=f\circ i_\lambda$; si $f$ est formellement lisse (resp. formellement non ramifié), il en est de même de $f\circ i_\lambda$ puisque $i_\lambda$ est formellement étale (17.1.3); mais comme $j_\lambda$ est formellement étale, cela entraîne que $f_\lambda$ est formellement lisse (resp. formellement non ramifié) en vertu de (17.1.5). Inversement, si tous les $f_\lambda$ sont formellement lisses (resp. formellement non ramifiés), il en est de même des $j_\lambda\circ f_\lambda$ (17.1.3), donc aussi de $f$ en vertu de (i).

